\documentclass[letterpaper]{article}

\usepackage[letterpaper,margin=1in]{geometry}

% Compile from the repository root so that this path resolves to
% the pinned WSP submodule.
\input{wsp/testing/test-case-library.tex}

\begin{document}

\def\TCID{%
  TC-0014
}

\def\TCTitle{%
  ListView items and columns
}

\def\TCPurpose{%
  Verify the ListView column headers, items, sub items, sorting, FindItem, hit
  testing, selection, errors and events.
}

\def\TCPriority{%
  \TCPriorityHigh
}

\def\TCRequirementRef{%
  REQ-0014, REQ-0032, REQ-0033, REQ-0034, REQ-0041
}

\def\TCDesignRef{%
  \texttt{ADR-0006} and \texttt{ADR-0010} in WFC
}

\def\TCFeatureRef{%
  \texttt{tests/corpus/192-comctl-listview.frm} and
  \texttt{tests/corpus/192-comctl-listview.vbp}; the form code and the
  controls of \texttt{include/wfc/mscomctl/mscomctl.inc} and the runner
  \texttt{tests/tc-0014-comctl-listview.ps1}
}

\def\TCTestTechnique{%
  Scenario-based testing of one form: each Immediate-window line is compared
  with the recorded outcome (ISO/IEC/IEEE 29119-4 experience-based and
  specification-based techniques, with expected output comparison).
}

\def\TCPreconditions{%
  A \texttt{wfc} of the same version as the library (WFC, the interpreter and
  compiler these controls belong to, runs the case: the library holds the
  program, WFC the runner). The files of the case are in
  \texttt{tests/corpus}: \texttt{192-comctl-listview.vbp},
  \texttt{192-comctl-listview.frm}, \texttt{192-comctl-listview.expected},
  \texttt{192-comctl-listview.debug}, \texttt{192-comctl-listview.ui}.
}

\def\TCEnvironment{%
  Windows with \texttt{wfc.exe} of the x64 Debug build; the form runs in WFC's
  headless host (no desktop is needed). WFC registers the case with CTest and
  runs \texttt{wfc --interpret 192-comctl-listview.vbp}.
}

\def\TCAssumptions{%
  The scripted input \texttt{192-comctl-listview.ui} is the only input; the
  case does not read the registry or other files except the picture named in
  the form.
}

\def\TCInputData{%
  A form with a ListView in Report view with check boxes; three columns
  (\texttt{name}, \texttt{size}, \texttt{kind}); items \texttt{Banana},
  \texttt{Apple}, \texttt{Carrot}; the script clicks, checks, clicks a column
  and edits.
}

\def\TCInitialState{%
  A fresh \texttt{wfc} process; the program creates its own controls from the
  form file and starts with no state from a previous run.
}

\def\TCProcedure{%
  \begin{enumerate}
    \item Add three columns and print the count, \texttt{Left} values,
      position, sub item index and alignment.
    \item Add items with sub items by \texttt{SubItems} and by
      \texttt{ListSubItems.Add} and print the counts, texts, tool tip and type
      name.
    \item Sort by the item, descending, and by a sub item.
    \item Print \texttt{FindItem} by text, partial text, missing text and sub
      item; \texttt{HitTest} and \texttt{GetFirstVisible}; select an item.
    \item Set \texttt{View} to 7, read a sub item past the end and add a
      duplicate key.
    \item Run the script.
  \end{enumerate}
}

\def\TCExpectedResult{%
  \begin{enumerate}
    \item Three columns with the \texttt{Left} values 0 and 2400 and the right
      alignment of \texttt{size} (1); three items and a sub item text
      \texttt{7}.
    \item Sorting gives \texttt{Apple Banana Carrot}, then \texttt{Carrot
      Apple} for the descending order and \texttt{Carrot Banana} for the sub
      item key.
    \item \texttt{FindItem} returns the right keys and \texttt{Nothing} for
      the unknown text.
    \item \texttt{View = 7} is 380, a sub item past the end is 35600 and a
      duplicate key is 35602.
    \item Events: \texttt{itemclick c}, \texttt{itemcheck cTrue},
      \texttt{column size}, \texttt{renamed Carrots}.
    \item \texttt{wfc} exits with code 0, writes the exact text of
      \texttt{192-comctl-listview.expected} to standard output, and writes the
      exact Immediate-window text of \texttt{192-comctl-listview.debug} to
      standard error. Any difference fails the case.
  \end{enumerate}
}

\def\TCPostConditions{%
  The form is closed by the last \texttt{close Form1} line of the script;
  nothing is left running and no file is written.
}

\TCTemplate

\end{document}
